
\documentclass{article}
\usepackage{colortbl}
\usepackage{makecell}
\usepackage{multirow}
\usepackage{supertabular}

\begin{document}

\newcounter{utterance}

\centering \large Interaction Transcript for game `hot\_air\_balloon', experiment `air\_balloon\_survival\_tr\_complexity\_hard', episode 2 with claude{-}opus{-}4{-}8{-}azure.
\vspace{24pt}

{ \footnotesize  \setcounter{utterance}{1}
\setlength{\tabcolsep}{0pt}
\begin{supertabular}{c@{$\;$}|p{.15\linewidth}@{}p{.15\linewidth}p{.15\linewidth}p{.15\linewidth}p{.15\linewidth}p{.15\linewidth}}
    \# & $\;$A & \multicolumn{4}{c}{Game Master} & $\;\:$B\\
    \hline

    \theutterance \stepcounter{utterance}  
    & & \multicolumn{4}{p{0.6\linewidth}}{
        \cellcolor[rgb]{0.9,0.9,0.9}{
            \makecell[{{p{\linewidth}}}]{
                \texttt{\tiny{[P1$\langle$GM]}}
                \texttt{İşbirliğine dayalı bir müzakere oyununa katılıyorsun.} \\
\\ 
\texttt{Başka bir katılımcıyla birlikte, elde tutulacak tek bir öğe kümesi üzerinde anlaşmanız gerekiyor. Her birinizin, her bir öğenin sizin için ne kadar önemli olduğuna dair kendi bakışı (önem değeri) vardır. Diğer katılımcının öğelere nasıl değer verdiğini bilmiyorsunuz. Ayrıca, her bir öğenin gerektirdiği çaba değeri de size verilmiştir.} \\
\texttt{Yalnızca seçilen öğelerin toplam çabası paylaşılan bir sınırı aşmıyorsa bir küme üzerinde anlaşabilirsiniz:} \\
\\ 
\texttt{LİMİT = 4843} \\
\\ 
\texttt{İşte öğelerin tek tek çaba değerleri:} \\
\\ 
\texttt{Öğe çabaları = \{"C08": 880, "C55": 158, "A43": 104, "A67": 514, "B42": 797, "A62": 815, "B19": 336, "C43": 79, "C73": 522, "C10": 973, "B36": 682, "A18": 178, "C61": 184, "C67": 795, "C42": 154, "C70": 145, "B40": 842, "C39": 887, "A21": 328, "C16": 313\}} \\
\\ 
\texttt{İşte her bir öğenin önemine dair kişisel bakışın:} \\
\\ 
\texttt{Öğe önem değerleri = \{"C08": 977, "C55": 255, "A43": 201, "A67": 611, "B42": 894, "A62": 912, "B19": 433, "C43": 176, "C73": 619, "C10": 1070, "B36": 779, "A18": 275, "C61": 281, "C67": 892, "C42": 251, "C70": 242, "B40": 939, "C39": 984, "A21": 425, "C16": 410\}} \\
\\ 
\texttt{Amaç:} \\
\\ 
\texttt{Amacın, LİMİT içinde kalarak sana mümkün olan en fazla faydayı sağlayan (yani SENİN için toplam önemi maksimize eden) ortak bir öğe kümesi üzerinde anlaşmaktır. Her mesajda mutlaka bir TEKLİF yapmak zorunda değilsin {-} sadece müzakere de edebilirsin. Tüm taktikler serbest!} \\
\\ 
\texttt{Etkileşim Protokolü:} \\
\\ 
\texttt{Bir mesajda yalnızca aşağıdaki yapılandırılmış biçimleri kullanabilirsin:} \\
\\ 
\texttt{TEKLİF: \{'A', 'B', 'C', …\}} \\
\texttt{Belirtilen öğelerin tutulmasını teklif eder.} \\
\texttt{REDDET: \{'A', 'B', 'C', …\}} \\
\texttt{Rakibin teklifini açıkça reddeder.} \\
\texttt{SAVUNMA: \{'...'\}} \\
\texttt{Son teklifinizi savunur veya diğer oyuncunun teklifine karşı argüman sunar.} \\
\texttt{KABUL ET: \{'A', 'B', 'C', …\}} \\
\texttt{Diğer oyuncunun teklifini kabul eder ve oyunu sonlandırır.} \\
\\ 
\\ 
\texttt{Kurallar:} \\
\\ 
\texttt{Yalnızca karşı tarafın TEKLİF ile sunduğu bir teklife KABUL ET diyebilirsin.} \\
\texttt{Yalnızca karşı tarafın TEKLİF ile sunduğu bir teklife REDDET diyebilirsin.} \\
\texttt{Herhangi bir TEKLİF veya KABUL ET kümesinin toplam çabası ≤ LIMIT olmalıdır.} \\
\texttt{Gizli önem değerlerini AÇIKLAMA.} \\
\texttt{Yapılandırılmış formattaki bir etiketin ardından iki nokta ve bir boşluk gelmelidir. İçerik, 0 veya daha fazla string içeren bir Python kümesi olmalıdır.} \\
\texttt{Yani format şu şekilde olmalıdır: "ETİKET: \{...\}"} \\
\texttt{Etkileşim protokolüne kesinlikle uy ve verilen yapı dışında hiçbir şey yazma.} \\
\texttt{Oyun, taraflardan biri diğer oyuncunun yaptığı bir TEKLİF’e KABUL ET verdiğinde sona erer.} \\
\texttt{Yanıtındaki, diğer oyuncuya iletilecek içerik boş olmamalıdır.} \\
\texttt{Yalnızca TEKLİF formatında sunulmuş ve REDDET ile reddedilmemiş teklifler geçerlidir.} \\
\\ 
            }
        }
    }
    & \\ \\

    \theutterance \stepcounter{utterance}  
    & \multicolumn{4}{p{0.6\linewidth}}{
        \cellcolor[rgb]{0.9,0.9,0.9}{
            \makecell[{{p{\linewidth}}}]{
                \texttt{\tiny{[P1$\rangle$GM]}}
                \texttt{TEKLİF: \{'C43', 'A43', 'C70', 'C42', 'C55', 'A18', 'C61', 'C16', 'A21', 'B19', 'A67', 'C73', 'B36', 'C10'\}} \\
            }
        }
    }
    & & \\ \\

    \theutterance \stepcounter{utterance}  
    & & \multicolumn{4}{p{0.6\linewidth}}{
        \cellcolor[rgb]{0.9,0.9,0.9}{
            \makecell[{{p{\linewidth}}}]{
                \texttt{\tiny{[GM$\rangle$P2]}}
                \texttt{İşbirliğine dayalı bir müzakere oyununa katılıyorsun.} \\
\\ 
\texttt{Başka bir katılımcıyla birlikte, elde tutulacak tek bir öğe kümesi üzerinde anlaşmanız gerekiyor. Her birinizin, her bir öğenin sizin için ne kadar önemli olduğuna dair kendi bakışı (önem değeri) vardır. Diğer katılımcının öğelere nasıl değer verdiğini bilmiyorsunuz. Ayrıca, her bir öğenin gerektirdiği çaba değeri de size verilmiştir.} \\
\texttt{Yalnızca seçilen öğelerin toplam çabası paylaşılan bir sınırı aşmıyorsa bir küme üzerinde anlaşabilirsiniz:} \\
\\ 
\texttt{LİMİT = 4843} \\
\\ 
\texttt{İşte öğelerin tek tek çaba değerleri:} \\
\\ 
\texttt{Öğe çabaları = \{"C08": 880, "C55": 158, "A43": 104, "A67": 514, "B42": 797, "A62": 815, "B19": 336, "C43": 79, "C73": 522, "C10": 973, "B36": 682, "A18": 178, "C61": 184, "C67": 795, "C42": 154, "C70": 145, "B40": 842, "C39": 887, "A21": 328, "C16": 313\}} \\
\\ 
\texttt{İşte her bir öğenin önemine dair kişisel bakışın:} \\
\\ 
\texttt{Öğe önem değerleri = \{"C08": 977, "C55": 255, "A43": 201, "A67": 611, "B42": 894, "A62": 912, "B19": 433, "C43": 176, "C73": 619, "C10": 1070, "B36": 779, "A18": 275, "C61": 281, "C67": 892, "C42": 251, "C70": 242, "B40": 939, "C39": 984, "A21": 425, "C16": 410\}} \\
\\ 
\texttt{Amaç:} \\
\\ 
\texttt{Amacın, LİMİT içinde kalarak sana mümkün olan en fazla faydayı sağlayan (yani SENİN için toplam önemi maksimize eden) ortak bir öğe kümesi üzerinde anlaşmaktır. Her mesajda mutlaka bir TEKLİF yapmak zorunda değilsin {-} sadece müzakere de edebilirsin. Tüm taktikler serbest!} \\
\\ 
\texttt{Etkileşim Protokolü:} \\
\\ 
\texttt{Bir mesajda yalnızca aşağıdaki yapılandırılmış biçimleri kullanabilirsin:} \\
\\ 
\texttt{TEKLİF: \{'A', 'B', 'C', …\}} \\
\texttt{Belirtilen öğelerin tutulmasını teklif eder.} \\
\texttt{REDDET: \{'A', 'B', 'C', …\}} \\
\texttt{Rakibin teklifini açıkça reddeder.} \\
\texttt{SAVUNMA: \{'...'\}} \\
\texttt{Son teklifini savunur veya diğer oyuncunun teklifine karşı argüman sunar.} \\
\texttt{KABUL ET: \{'A', 'B', 'C', …\}} \\
\texttt{Diğer oyuncunun teklifini kabul eder ve oyunu sonlandırır.} \\
\\ 
\\ 
\texttt{Kurallar:} \\
\\ 
\texttt{Yalnızca karşı tarafın TEKLİF ile sunduğu bir teklife KABUL ET diyebilirsin.} \\
\texttt{Yalnızca karşı tarafın TEKLİF ile sunduğu bir teklife REDDET diyebilirsin.} \\
\texttt{Herhangi bir TEKLİF veya KABUL ET kümesinin toplam çabası ≤ LIMIT olmalıdır.} \\
\texttt{Gizli önem değerlerini AÇIKLAMA.} \\
\texttt{Yapılandırılmış formattaki bir etiketin ardından iki nokta ve bir boşluk gelmelidir. İçerik, 0 veya daha fazla string içeren bir Python kümesi olmalıdır.} \\
\texttt{Yani format şu şekilde olmalıdır: "ETİKET: \{...\}"} \\
\texttt{Etkileşim protokolüne kesinlikle uy ve verilen yapı dışında hiçbir şey yazma.} \\
\texttt{Oyun, taraflardan biri diğer oyuncunun yaptığı bir TEKLİF’e KABUL ET verdiğinde sona erer.} \\
\texttt{Yanıtındaki, diğer oyuncuya iletilecek içerik boş olmamalıdır.} \\
\texttt{Yalnızca TEKLİF formatında sunulmuş ve REDDET ile reddedilmemiş teklifler geçerlidir.} \\
\\ 
\\ 
\\ 
\texttt{Şimdi diğer oyuncunun ilk mesajını alacaksın.} \\
\\ 
\texttt{TEKLİF: \{'C43', 'A43', 'C70', 'C42', 'C55', 'A18', 'C61', 'C16', 'A21', 'B19', 'A67', 'C73', 'B36', 'C10'\}} \\
            }
        }
    }
    & \\ \\

    \theutterance \stepcounter{utterance}  
    & & & \multicolumn{4}{p{0.6\linewidth}}{
        \cellcolor[rgb]{0.9,0.9,0.9}{
            \makecell[{{p{\linewidth}}}]{
                \texttt{\tiny{[GM$\langle$P2]}}
                \texttt{KABUL ET: \{'C43', 'A43', 'C70', 'C42', 'C55', 'A18', 'C61', 'C16', 'A21', 'B19', 'A67', 'C73', 'B36', 'C10'\}} \\
            }
        }
    }
     \\ \\

    \theutterance \stepcounter{utterance}  
    & & & \multicolumn{2}{p{0.3\linewidth}}{
        \cellcolor[rgb]{0.9,0.9,0.9}{
            \makecell[{{p{\linewidth}}}]{
                \texttt{\tiny{[GM$|$GM]}}
                \texttt{game successful} \\
            }
        }
    }
    & & \\ \\

    \theutterance \stepcounter{utterance}  
    & & & \multicolumn{2}{p{0.3\linewidth}}{
        \cellcolor[rgb]{0.9,0.9,0.9}{
            \makecell[{{p{\linewidth}}}]{
                \texttt{\tiny{[GM$|$GM]}}
                \texttt{end game} \\
            }
        }
    }
    & & \\ \\

\end{supertabular}
}

\end{document}
